\documentclass[10pt,a4paper]{amsart}
\usepackage[margin=19mm]{geometry}
\usepackage{amsmath,amssymb,mathtools}
\usepackage[scr=rsfs,cal=euler]{mathalfa}
\usepackage{xfrac}
\usepackage[unicode,hidelinks,bookmarks=false]{hyperref}
\newcommand{\ie}{i.e.\ }
\newcommand{\cf}{cf.\ }
\newcommand{\resp}{respectively\ }
\newcommand{\Subsection}{\subsection}
\providecommand{\nobreakdash}{\nobreak\hbox{-}}
\setlength{\emergencystretch}{2em}
\multlinegap0pt
\usepackage{bm}
\usepackage{bbm}
\usepackage{dsfont}
\usepackage{xspace}
\usepackage{cancel}
\usepackage{mathtools}
\usepackage{amsthm}
\makeatletter
\@ifundefined{Theorem}{\newtheorem{Theorem}{Theorem}}{}
\@ifundefined{Lemma}{\newtheorem{Lemma}[Theorem]{Lemma}}{}
\makeatother
\newcommand{\Hom}{\operatorname{Hom}}
\newcommand{\Res}{\operatorname{Res}}
\newcommand{\bZ}{\mathbb{Z}}
\newcommand{\bv}{\mathbf{v}}
\newcommand{\cH}{\mathcal{H}}
\newcommand{\id}{\operatorname{id}}
\newcommand{\str}{\operatorname{Str}}
\title[Shifted quantum groups via critical stable envelopes]{Shifted quantum groups via critical stable envelopes}
\author{Yalong Cao, Andrei Okounkov, Yehao Zhou, Zijun Zhou}
\date{Source: arXiv:2601.01518v1 (CC BY 4.0)}
\begin{document}
\maketitle

\noindent\emph{Source and licence.} Shifted quantum groups via critical stable envelopes. Version arXiv:2601.01518v1 (CC BY 4.0), distributed under the Creative Commons Attribution 4.0 International licence, as recorded in the arXiv licence metadata for this exact version and shown on the abstract page https://arxiv.org/abs/2601.01518v1. Full rights notice: licenses/arxiv-2601.01518-CC-BY-4.0.txt.\par\smallskip

\noindent\textit{Editorial scope.} Counted unit: the Lemma whose source label is \texttt{lem fund induction} in the article (source0.tex lines 2458--2465, source label \texttt{lem fund induction}), delivered together with the article's own complete original proof of it (source0.tex lines 2467--2507, one \texttt{proof} environment of 41 lines). No mathematics is written, repaired or re-derived: statement and proof are the article's own text line for line. Required context retained uncounted: none. Every definition, display and label of those lines is retained unchanged. Omitted from this package and not redistributed: 1--2457, 2466--2466, 2508--8514. Nothing inside a retained excerpt is omitted or altered: every retained body is delivered exactly as the article's lines are written, so no omission receipt of any kind accompanies this package. Citations: the retained excerpts carry no citation command at all, so no bibliography entry of the article is transcribed and no reference is redistributed. Reference adaptation: none was needed and none was made; every reference inside the delivered unit resolves inside it, so no unresolved reference or citation is produced. Printed slips kept literal: no printed word, symbol, hypothesis or number of the counted statement or of its proof was altered. Layout and class: the article's own class is \texttt{amsart} with packages including color, graphicx, tikz, xcolor, cancel, amssymb, amsfonts, amsthm, amsmath, calligra, enumerate, extarrows, inputenc, accents; the wrapper is the lane's pinned \texttt{amsart} editorial wrapper at 10pt on A4, which restates the theorem-like environments the retained text needs and adds the article's own macro definitions that the retained text needs (\texttt{\textbackslash Hom}, \texttt{\textbackslash Res}, \texttt{\textbackslash bZ}, \texttt{\textbackslash bv}, \texttt{\textbackslash cH}, \texttt{\textbackslash id}, \texttt{\textbackslash str}), with extra wrapper packages (bm, bbm, dsfont, xspace, cancel, mathtools). The delivered numbering is the unit's own. Assets: no retained span contains a figure environment, a graphics-inclusion command, a PDF-page inclusion, a table or a TikZ picture; no image file is copied into this package, the package declares no asset dependency and no image file is redistributed with it. Licence: version arXiv:2601.01518v1 of this article is distributed under the Creative Commons Attribution 4.0 International licence, as recorded in the arXiv licence metadata for this exact version and shown on the abstract page https://arxiv.org/abs/2601.01518v1. A full rights notice is delivered at licenses/arxiv-2601.01518-CC-BY-4.0.txt. Characters: every retained excerpt is pure ASCII, so the delivered engine, XeLaTeX with its default Latin Modern text font, reports no missing character.
\begin{Lemma}\label{lem fund induction}
Let $\mathfrak{c}_1,\mathfrak{c}_2$ be arbitrary auxiliary data, $|\mathbf 0\rangle$ be the vacuum of $ \cH_{\mathfrak{c}_1}(\mathbf 0)$, $m\in \Hom(\cH_{\mathfrak{c}_2}(\bv_1),\cH_{\mathfrak{c}_2}(\bv_2))$. Then for any $k\in \bZ_{\geqslant 0}$, we have
\begin{align}\label{fund induction}
[\mathsf E(|\mathbf 0\rangle\langle\mathbf 0|\: u),\mathsf E(m\: u^{k})]-[\mathsf E(|\mathbf 0\rangle\langle\mathbf 0|),\mathsf E(m\: u^{k+1})]&=\mathsf E(m'\: u^{k})+\mathsf E(m''\: u^{k})\:\mathsf E(|\mathbf 0\rangle\langle\mathbf 0|)\\ \nonumber
&\quad +\sum_{\bv_2\geqslant\bv> 0}\mathsf E(f_{\bv}\: u^{k})\mathsf E(g_{\bv})+\sum_{\bv_1\geqslant\bv> 0}\mathsf E(f'_{\bv}\: u^{k}) \mathsf E(g'_{\bv}),
\end{align}
for certain $m',m''\in \Hom(\cH_{\mathfrak{c}_2}(\bv_1),\cH_{\mathfrak{c}_2}(\bv_2))$, $g_{\bv}\in \Hom(\cH_{\mathfrak{c}_1}(\mathbf 0),\cH_{\mathfrak{c}_1}(\bv))$, $f_{\bv}\in \Hom(\cH_{\mathfrak{c}_2}(\bv_1),\cH_{\mathfrak{c}_2}(\bv_2-\bv))$, $g'_{\bv}\in \Hom(\cH_{\mathfrak{c}_1}(\bv),\cH_{\mathfrak{c}_1}(\mathbf 0))$, $f'_{\bv}\in \Hom(\cH_{\mathfrak{c}_2}(\bv_1-\bv),\cH_{\mathfrak{c}_2}(\bv_2))$.
\end{Lemma}

\begin{proof}
By the RTT=TTR equation, we have
\begin{align*}
[\mathsf E(|\mathbf 0\rangle\langle\mathbf 0|\: v),\mathsf E(m\: u^{k})]=\underset{v\to \infty}{\Res}\underset{u\to \infty}{\Res}\str_{\cH_{\mathfrak{c}_1}\otimes \cH_{\mathfrak{c}_2}}\left[\left(R_{12}(v-u)(|\mathbf 0\rangle\langle\mathbf 0|\: v\:\otimes m\: u^{k})R_{12}(v-u)^{-1}- |\mathbf 0\rangle\langle\mathbf 0|\: v\:\otimes m\: u^{k}\right)T_2(u)T_{1}(v)\right],
\end{align*}
where $T_i$ $(i=1,2)$ is the $R$-matrix that braids $\cH_{\mathfrak{c}_i}$ with a state space. By taking expansions 
$$R_{12}(v-u)=\id+\frac{\pmb r_{12}}{v-u}+\frac{\pmb s}{(v-u)^2}+O((v-u)^{-3}),$$ 
$$R_{12}(v-u)^{-1}=\id-\frac{\pmb r_{12}}{v-u}+\frac{\pmb s'}{(v-u)^2}+O((v-u)^{-3}),$$ 
we get
\begin{multline*}
[\mathsf E(|\mathbf 0\rangle\langle\mathbf 0|\: v),\mathsf E(m\: u^{k})]=\underset{v\to \infty}{\Res}\underset{u\to \infty}{\Res}\str_{\cH_{\mathfrak{c}_1}\otimes \cH_{\mathfrak{c}_2}}\\
\left[\left([\pmb r_{12},|\mathbf 0\rangle\langle\mathbf 0| \otimes m](u^k+v^{-1}u^{k+1})+\pmb s\cdot(|\mathbf 0\rangle\langle\mathbf 0| \otimes m)v^{-1}u^k+(|\mathbf 0\rangle\langle\mathbf 0| \otimes m)\cdot \pmb s'\:v^{-1}u^k\right)T_2(u)T_{1}(v)\right].
\end{multline*}
We ignore the higher order in $u/v$ terms in the above expansion because $T_1(v)$ has no pole as $v\to \infty$. The same argument shows that
\begin{align*}
[\mathsf E(|\mathbf 0\rangle\langle\mathbf 0|\: v),\mathsf E(m\: u^{k})]=\underset{v\to \infty}{\Res}\underset{u\to \infty}{\Res}\str_{\cH_{\mathfrak{c}_1}\otimes \cH_{\mathfrak{c}_2}}
\left[[\pmb r_{12},|\mathbf 0\rangle\langle\mathbf 0| \otimes m]v^{-1}u^{k+1}\:T_2(u)T_{1}(v)\right].
\end{align*}
It follows that
\begin{multline*}
\text{LHS of \eqref{fund induction}}=\underset{v\to \infty}{\Res}\underset{u\to \infty}{\Res}\str_{\cH_{\mathfrak{c}_1}\otimes \cH_{\mathfrak{c}_2}}
\\
\left[\left([\pmb r_{12},|\mathbf 0\rangle\langle\mathbf 0| \otimes m]\:u^{k}+\pmb s\cdot(|\mathbf 0\rangle\langle\mathbf 0| \otimes m)v^{-1}u^k+(|\mathbf 0\rangle\langle\mathbf 0| \otimes m)\cdot \pmb s'\:v^{-1}u^k\right)T_2(u)T_{1}(v)\right].
\end{multline*}
Since the nonvanishing terms in $\lim_{v\to \infty}T_1(v)$ are concentrated in the diagonal, we have 
\begin{align*}
    \underset{v\to \infty}{\Res}\underset{u\to \infty}{\Res}\str_{\cH_{\mathfrak{c}_1}\otimes \cH_{\mathfrak{c}_2}}
\left[\left(\pmb s\cdot(|\mathbf 0\rangle\langle\mathbf 0| \otimes m)v^{-1}u^k+(|\mathbf 0\rangle\langle\mathbf 0| \otimes m)\cdot \pmb s'\:v^{-1}u^k\right)T_2(u)T_{1}(v)\right]=\mathsf E\left(m'\:u^k \right),
\end{align*}
where $m'=\langle\mathbf 0|\pmb s |\mathbf 0\rangle \cdot m+m\cdot\langle\mathbf 0|\pmb s' |\mathbf 0\rangle$ is a homomorphism from $\cH_{\mathfrak{c}_2}(\bv_1)$ to $\cH_{\mathfrak{c}_2}(\bv_2)$. On the other hand, we can write the other term in the LHS of \eqref{fund induction} in homogeneous components and get 
\begin{align*}
    \underset{v\to \infty}{\Res}\underset{u\to \infty}{\Res}\str_{\cH_{\mathfrak{c}_1}\otimes \cH_{\mathfrak{c}_2}}
\left[\pmb r_{12}\:(|\mathbf 0\rangle\langle\mathbf 0| \otimes m)u^{k}\:T_2(u)T_{1}(v)\right]=\sum_{\bv_2\geqslant\bv\geqslant 0}\mathsf E(f_{\bv}\: u^{k})\mathsf E(g_{\bv}),
\end{align*}
where $g_{\bv}\otimes f_{\bv}$ is the short-hand notation for the $\Hom(\cH_{\mathfrak{c}_1}(\mathbf 0)\otimes \cH_{\mathfrak{c}_2}(\bv_1),\cH_{\mathfrak{c}_1}(\bv)\otimes \cH_{\mathfrak{c}_2}(\bv_2-\bv))$ component of $\pmb r_{12}\:(|\mathbf 0\rangle\langle\mathbf 0| \otimes m)$. $g_{\bv}\otimes f_{\bv}$ should be actually a linear combination of this form but we suppress the summation for simplicity. Similarly,
\begin{align*}
    \underset{v\to \infty}{\Res}\underset{u\to \infty}{\Res}\str_{\cH_{\mathfrak{c}_1}\otimes \cH_{\mathfrak{c}_2}}
\left[(|\mathbf 0\rangle\langle\mathbf 0| \otimes m)\:\pmb r_{12}\:u^{k}\:T_2(u)T_{1}(v)\right]=\sum_{\bv_2\geqslant\bv\geqslant 0}\mathsf E(f'_{\bv}\: u^{k})\mathsf E(g'_{\bv})
\end{align*}
for $g'_{\bv}\otimes f'_{\bv}$ being the $\Hom(\cH_{\mathfrak{c}_1}(\bv)\otimes \cH_{\mathfrak{c}_2}(\bv_1-\bv),\cH_{\mathfrak{c}_1}(\mathbf 0)\otimes \cH_{\mathfrak{c}_2}(\bv_2))$ component of $(|\mathbf 0\rangle\langle\mathbf 0| \otimes m)\:\pmb r_{12}$. Combine the above three equations and let $\mathsf E(m''\: u^{k})\:\mathsf E(|\mathbf 0\rangle\langle\mathbf 0|)=\mathsf E(f_{\mathbf 0}\: u^{k})\mathsf E(g_{\mathbf 0})+\mathsf E(f'_{\mathbf 0}\: u^{k})\mathsf E(g'_{\mathbf 0})$, we get \eqref{fund induction}.
\end{proof}

\end{document}
